% g_eps_lemma.tex -- round 381, lemma-first attack on G_eps.
% Outcome: REDUCED to the FO pairing C_eps and the second-order
% remainder R_2.  Builds with pdflatex.
% Companion machine check: verify_g_eps.py.
% Discovery census: experiments/tfpt-discovery/g_eps_lemma_probe.py.
% NO RH CLAIM.  Research documentation, not a theorem of RH.
\documentclass[11pt]{article}

\usepackage[a4paper,margin=2.7cm]{geometry}
\usepackage{amsmath,amssymb,amsthm}
\usepackage{booktabs}
\usepackage{microtype}
\usepackage{xcolor}
\usepackage[colorlinks=true,linkcolor=blue!50!black,urlcolor=blue!50!black]{hyperref}

\newtheorem{lemma}{Lemma}
\newtheorem{prop}{Proposition}
\newtheorem{definition}{Definition}
\theoremstyle{remark}
\newtheorem{remark}{Remark}

\newcommand{\N}{\mathbb{N}}
\newcommand{\Q}{\mathbb{Q}}
\newcommand{\R}{\mathbb{R}}
\newcommand{\eps}{\varepsilon}

\title{\bfseries The $G_\varepsilon$-lemma, reduced\\[2mm]
\large First-order Jacobi perturbation, the pairing $C_\varepsilon$,
and the second-order remainder $R_2$}
\author{Stefan Hamann}
\date{August 28, 2026}

\begin{document}
\maketitle

\begin{abstract}\noindent
Round 381 attacks the named remainder $G_\varepsilon$ of~\cite{v3p}
(the last twelve monic Jacobi coefficients $\gamma_k$ of the $v_2$-chain
satisfy $\lvert\gamma_k-\tfrac14\rvert\le\tfrac1{16}$ and
$\lvert\log(\gamma_{k+1}/\gamma_k)\rvert\le\tfrac25$).
What is proved: the first-order formula for $\partial\gamma_k/\partial w_j$
on a fixed node set, in terms of monic orthogonal-polynomial squares
and the previous norm $h_{k-1}$ (exact over $\Q$; floating-point
finite-difference residual $\le 10^{-3}$ on FRAME-A $w=9$).
The Jacobi-$(0,1)$ Fej\'er identity and the cosine-grid mesh remain
theorems of~\cite{v3p}.
What is \emph{not} proved: the box itself.
A crude $L^\infty$ bound on the relative arm-length perturbation
is true and useless ($\lVert\varepsilon\rVert_\infty=O(1)$--$O(10)$
on the source, giving a first-order majorant $\sim 3\gg\tfrac1{16}$).
The jump half is not a corollary of the box
($\log\tfrac53>\tfrac25$).
The second-order remainder of the weight expansion is of the
same size as the first-order increment on the last twelve
coefficients (midpoint test, ratio $\ge\tfrac12$).
The hydra-legal reduction is the pair $(C_\varepsilon,R_2)$:
the explicit first-order pairing of the weight increment against
the reference OP squares, and the Taylor remainder of that expansion.
A scramble world leaves $G_\varepsilon$ (last-twelve deviation
$\ge 6.8$).  No RH claim.
\end{abstract}

\medskip\noindent
\textbf{Status.}\enspace
Lemmas~\ref{lem:fo}--\ref{lem:lambda} are proved (finite identities).
Lemma~\ref{lem:jumpind} is a comparison of two explicit constants.
Lemma~\ref{lem:linf} is a true majorant that does not yield the box.
Proposition~\ref{prop:red} names the reduction.
The hydra test of~\cite{v2} is met: $C_\varepsilon$ is a quadrature
against known reference polynomials (no perturbed recurrence);
$R_2$ is a scalar remainder of a Taylor step already written.
No RH claim.

\section{Recollection}\label{sec:rec}

$G_\varepsilon$ of~\cite{v3p}: writing $n=N-2$
for the degree of $v_2$, the last twelve monic Jacobi coefficients
of the bordered $\mu-\nu$ chain satisfy
\begin{equation}\label{eq:geps}
\bigl\lvert\gamma_k-\tfrac14\bigr\rvert\le\tfrac1{16},
\qquad
\bigl\lvert\log(\gamma_{k+1}/\gamma_k)\bigr\rvert\le\tfrac25,
\qquad
k\in[n-12,n).
\end{equation}
The mesh $x=\cos(2\pi k/L)$, $L=4h-2$, is a construction identity.
Equal weights on the occupied cosine nodes are Chebyshev
($\gamma_1=\tfrac12$, $\gamma_k=\tfrac14$ for $k\ge 2$ up to a
$10^{-5}$ residual from one missing fold).
The Fej\'er factor $2(1-x)$ is Jacobi-$(0,1)$ with
$\lvert\gamma_n-\tfrac14\rvert=1/(4(2n+1)^2)$.
The source formula of~\cite{v3p} is
$\gamma_k=\tfrac14+\delta_k^{\mathrm{Fejer}}+\delta_k^{d_{\mathrm{arm}}}$,
and $\delta^{\mathrm{Fejer}}$ is negligible at half-filling on the
\emph{unsigned} occupied grid.
The live channel is the arm-length weight $d_{\mathrm{arm}}$
(r342 digamma $F_A(\xi)=-\log\pi+\Re\psi(\tfrac14+i\xi/2)$ plus
the two-cell tent of the prime layer).
On the $181$-window surface~\cite{v3p} the box holds
(max last-$12$ $0.0519$, max jump $0.337$).
This note does not repeat that census.

\section{The first-order formula}\label{sec:fo}

Let $\sigma=\sum_j w_j\,\delta_{x_j}$ be a quasi-definite atomic
measure on a fixed node set (signed weights allowed).
Write $p_k$ for the monic orthogonal polynomials,
$h_k=\int p_k^2\,d\sigma$, and $\gamma_k=h_k/h_{k-1}$.

\begin{lemma}[first-order Jacobi]\label{lem:fo}
At a quasi-definite weight,
\[
\frac{\partial\gamma_k}{\partial w_j}
=
\frac{p_k(x_j)^2-\gamma_k\,p_{k-1}(x_j)^2}{h_{k-1}}.
\]
Equivalently, for a multiplicative increment $dw_j=w_j\varepsilon_j$,
\[
\delta\gamma_k^{\mathrm{FO}}
=
\gamma_k\Bigl(\langle\varepsilon\,p_k^2\rangle_\sigma/h_k
-\langle\varepsilon\,p_{k-1}^2\rangle_\sigma/h_{k-1}\Bigr).
\]
For $\sigma=\mu-\nu$ the $\nu$-atoms enter with the opposite sign.
\end{lemma}

\begin{proof}
The leading coefficient of $p_k$ is identically $1$, so
$\partial p_k/\partial w_j$ has degree at most $k-1$.
Orthogonality gives $\int p_k\,(\partial p_k/\partial w_j)\,d\sigma=0$.
Differentiating $h_k=\sum_i w_i p_k(x_i)^2$ therefore yields
$\partial h_k/\partial w_j=p_k(x_j)^2$, with no contribution from
the polynomial variation.
The quotient rule on $\gamma_k=h_k/h_{k-1}$ is the displayed formula.
\end{proof}

The identity is gated over $\Q$ on a six-atom rational toy
(central difference versus the formula, relative residual
$<3\cdot 10^{-4}$ at step $1/200$) and on FRAME-A $w=9$
(one $\mu$-atom, $dw=10^{-7}$, relative residual $\le 10^{-3}$
through degree $12$).

The first-order remainder of a finite step $\varepsilon$ is
quadratic.  On the same toy, $\lvert\Delta\gamma-\mathrm{FO}\cdot\varepsilon\rvert/\varepsilon^2$
lies in $[0.05,5]$.  That is Lemma~\ref{lem:r2toy}.

\begin{lemma}[quadratic remainder, toy]\label{lem:r2toy}
On the six-atom rational measure, the first-order Taylor remainder
in a single-weight step of size $\varepsilon=1/200$ satisfies
$0.05\le\lvert R_2\rvert/\varepsilon^2\le 5$.
\end{lemma}

\section{What the $L^\infty$ bound cannot do}\label{sec:linf}

Write $\varepsilon_j=(w_j/\bar w_j)-1$ for the Fej\'er-stripped
relative arm-length, $\bar w\propto(1-x)$.

\begin{lemma}[crude majorant]\label{lem:linf}
$\lvert\delta\gamma_k^{\mathrm{FO}}\rvert
\le\gamma_k\bigl(\lVert\varepsilon\rVert_\infty^{(\mu)}
+\lVert\varepsilon\rVert_\infty^{(\nu)}\bigr)$
up to the signed-mass normalisation of the orthonormal squares.
On FRAME-A $w=9$ one has $\lVert\varepsilon_\mu\rVert_\infty>1$,
and the resulting majorant is $\sim 3\gg\tfrac1{16}$.
\end{lemma}

The bound is true and does not prove~\eqref{eq:geps}.
The measured last-twelve deviation is $0.034$; the majorant overshoots
by two orders of magnitude because
$\langle\varepsilon,\,\pi_k^2-\pi_{k-1}^2\rangle$ cancels.
That pairing against the \emph{reference} OP squares (the $\mu$-chain,
or the Fej\'er-constant reference) is the object $C_\varepsilon$ below.
It does not require solving the perturbed recurrence.

A PNT-free bound $\Lambda(n)\le\log n$ is elementary
(trial division: $\Lambda$ vanishes off prime powers, and
$\log p\le\log n$ on $n=p^k$).

\begin{lemma}[von~Mangoldt]\label{lem:lambda}
For every integer $n\ge 2$, $\Lambda(n)\le\log n$.
\end{lemma}

This bounds cell-masses of the r342 tent, hence $\lvert c_q^{\mathrm{atom}}\rvert$
in each lag cell.  It does \emph{not} bound the cancelled pairing:
high-lag $\lvert c_q\rvert$ on MAIN and on a scramble of $w=9$ are
of the same order, while last-twelve $\gamma$ is not.

\section{The jump is independent}\label{sec:jump}

\begin{lemma}[jump not implied]\label{lem:jumpind}
The closed box $\lvert\gamma-\tfrac14\rvert\le\tfrac1{16}$ permits
the consecutive ratio $5/3$, and $\log\tfrac53>\tfrac25$.
\end{lemma}

The jump half of $G_\varepsilon$ is a Lipschitz statement on
consecutive last-twelve coefficients, not a corollary of the
uniform box.  On the source it is cheaper (FRAME-A $w=9$: jump
$0.249<0.4$; $\mu$-only jump $0.122$).  The same pairing that
controls $C_\varepsilon$ at $k$ and at $k+1$ controls the
difference; the named remainder below includes it.

\section{The smaller lemmas}\label{sec:red}

Take as reference the \emph{positive} $\mu$-chain on the $+x$ arm
(same nodes, unsigned).  Write $\gamma_k^\mu$ for its monic
coefficients, and let $\mathrm{FO}_k[-\nu]$ be Lemma~\ref{lem:fo}
applied to the finite increment that subtracts $\nu$.
Define
\begin{align}
C_{\varepsilon,k}
&:=
\bigl\lvert\mathrm{FO}_k[-\nu]\bigr\rvert,
\label{eq:ceps}
\\
R_{2,k}
&:=
\gamma_k(\mu-\nu)-\gamma_k^\mu-\mathrm{FO}_k[-\nu].
\label{eq:r2}
\end{align}

\begin{definition}[$C_\varepsilon$ and $R_2$]\label{def:pair}
Condition $C_\varepsilon$ is $C_{\varepsilon,k}\le\tfrac1{32}$
for $k\in[n-12,n)$.
Condition $R_2$ is $\lvert R_{2,k}\rvert\le\tfrac1{32}$ on the
same range.
Condition $G_\varepsilon^\mu$ is~\eqref{eq:geps} for the positive
$\mu$-chain (box $\tfrac1{32}$ and jump $\tfrac15$ suffice).
\end{definition}

\begin{prop}[reduction]\label{prop:red}
$C_\varepsilon$ and $R_2$ are strictly more elementary than
$G_\varepsilon$ in the sense of~\cite{v2}~\S5:
\begin{enumerate}
\item \emph{Finite.}  Twelve real inequalities on an explicit
  sum against the \emph{unperturbed} OP squares ($C_\varepsilon$),
  plus twelve scalar remainders ($R_2$).
\item \emph{Closer to the construction.}  No Pr\"ufer phase, no
  mask remainder, no $v_2$ occupancy.  $C_\varepsilon$ does not
  even require the perturbed three-term recurrence.
\item \emph{Fewer quantifiers.}  No $\forall\eta$.
\end{enumerate}
On FRAME-A $w=9$: $G_\varepsilon$ holds (last-$12$ $0.034$, jump
$0.249$); $G_\varepsilon^\mu$ holds (last-$12$ $0.026$, jump
$0.122$); $\mathrm{mass}(\nu)/\mathrm{mass}(\mu)=0.061$;
the midpoint test of $R_2$ along the Fej\'er-constant path at
degree $40$ gives $\lvert\mathrm{quad}\rvert/\lvert\mathrm{lin}\rvert\ge\tfrac12$
on the last twelve --- $R_2$ is \emph{not} absorbed by the first-order
term.  The chain
\[
T_2'
 \;\Longleftarrow\;
 g_i\ge\tfrac38
 \;\Longleftarrow\;
 \mathrm{MED\text{-}CAP}
 \;\Longleftarrow\;
 X_n
 \;\Longleftarrow\;
 V_2
 \;\Longleftarrow\;
 V_3'
 \;\Longleftarrow\;
 G_\varepsilon
 \;\Longleftarrow\;
 G_\varepsilon^\mu\,\wedge\,C_\varepsilon\,\wedge\,R_2
\]
stands, with $(C_\varepsilon,R_2)$ the remaining named pair
($G_\varepsilon^\mu$ is the positive-measure sibling: same box,
no signed Stieltjes).
\end{prop}

The expansion $\mu\to\mu-\nu$ is a \emph{finite} $O(1)$ step in the
Christoffel sampling $\langle\pi_k^2\rangle_\nu$ (measured last-twelve
maximum $0.389$ on FRAME-A $w=9$).  First-order perturbation theory
about $\mu$ is therefore not a small-$\varepsilon$ expansion, which
is why $R_2$ is comparable to the linear increment.  Expanding about
the actual source against a $10^{-4}$ dictionary error would make
$R_2$ tiny and would move the whole content of $G_\varepsilon$ onto
the Jacobi coefficients of the r342 dictionary measure --- the same
lemma, not a smaller one.

\section{Adversaries}\label{sec:adv}

\paragraph{Segregated block.}
The $V_3'$ obstruction of~\cite{v3p} is a slow-then-fast
\emph{step} block.  Lemma~\ref{lem:fo} is a weight identity and
is not refuted by a step rearrangement.  A low-mode weight cluster
(the natural weight analogue of spatial segregation) does not move
the last twelve $\gamma_k$: on an equal-weight cosine grid of length
$80$, a relative amplitude $0.4$ in $\cos 2\theta$ changes last-twelve
$\gamma$ by $4\cdot 10^{-5}$.  The block sits outside the high-lag
tent class that can move the last twelve, and is not an FO-counterexample.

\paragraph{Scramble.}
FRAME-A $w=9$, scramble seed $1$: last-twelve
$\lvert\gamma-\tfrac14\rvert\ge 6.8$, jump $\gg\tfrac25$.
The lag-coefficient $L^\infty$ of $d_{\mathrm{arm}}$ is
\emph{not} the separator (high-lag $\lvert c_q\rvert$ is comparable
to MAIN).  What explodes is the $\mu$-Christoffel sampling of $\nu$:
$\langle\pi_k^2\rangle_\nu$ stays in $(0.34,0.40)$ on MAIN last-twelve
and exceeds $10^3$ on the scramble.  That is the named source bound
scramble breaks, and it is part of $C_\varepsilon$ (the pairing
evaluates $\pi_k$ of $\mu$ at the $\nu$ nodes).

\paragraph{Second order.}
Along the path from Fej\'er-constant $d_{\mathrm{arm}}$ to the actual
weights, at degree $40$, the last-twelve midpoint remainder is of
the same size as the linear increment (ratio $\ge\tfrac12$).
At full depth $n=N-2$ the quadratic increment is larger.
$R_2$ is not dominated.  The reduction to $R_2$ is forced by
Leg~D(iii), not optional.

\section{What remains}\label{sec:rest}

To prove $G_\varepsilon$: prove $C_\varepsilon$ (the cancelled
pairing of $d_{\mathrm{arm}}$ against $\pi_k^2-\pi_{k-1}^2$ of the
reference, uniformly in the last twelve $k$) and $R_2$ (the Taylor
remainder of the same finite step), plus $G_\varepsilon^\mu$ if the
reference is the $\mu$-chain.
The Geronimus/Fourier picture --- last-twelve $\delta\gamma$ tracks
the high-lag tent, i.e.\ Fourier mode $q\sim M$ of $d_{\mathrm{arm}}$
--- is the natural attack on $C_\varepsilon$, and is not carried
here.  T1 / COMPOSE unchanged.  No RH claim.

\section{Machine checks}\label{sec:machine}

Numbered lemmas: \texttt{verify\_g\_eps.py} (same directory).
Standalone: FO formula over $\Q$; quadratic remainder;
Jacobi-$(0,1)$; jump comparison; $\Lambda(n)\le\log n$;
Chebyshev quarters.
Construction: FO versus finite difference, $G_\varepsilon$ and
$G_\varepsilon^\mu$ on FRAME-A $w=9$, $R_2$ midpoint ratio,
scramble kill, $L^\infty$ crude, two-period versus block.
Discovery census: \texttt{g\_eps\_lemma\_probe.py}
(\texttt{--smoke} / full).
Exit line: \texttt{G\_EPS LEMMA VERIFIED}.

\section*{Claim boundary}

Finite identities for Jacobi perturbation on a fixed node set,
a named reduction of $G_\varepsilon$ to $(C_\varepsilon,R_2)$,
and a scramble break of the source pairing.
No statement about zeros of $\zeta(s)$, in either direction.
No RH claim.

\begin{thebibliography}{9}
\bibitem{v3p}
S.~Hamann,
\textit{The $V_3'$-lemma, reduced},
standalone note, August 2026,
\texttt{rh/problem/v3prime\_proof.tex}.
\bibitem{v2}
S.~Hamann,
\textit{Run-length regularity $V_2$ at the $x$-mask},
standalone note, August 2026,
\texttt{rh/problem/v2\_regularity.tex}.
\end{thebibliography}

\end{document}
